% ==========================================================================
%  Chapter 107 -- Macro-Guided Kinematic Regularization of Ko-Bathe Reversals
% ==========================================================================

\chapter{Macro-Guided Kinematic Regularization of Ko--Bathe Reversals}
\label{ch:kinematic-regularization}

\paragraph{Normative status.}
Third remediation round, building on the FE\textsuperscript{2} probe of
Cap.~\ref{ch:energy-linesearch-fe2}. The working hypothesis, stated by
the author and supported by the two-way evidence, is that the macro
model \emph{guides} the local model toward the physical branch because
imposed section kinematics are a strong control, and that enriching the
kinematic transfer (interior sections, and eventually the full
interpolated field) regularizes the Ko--Bathe reversal behaviour at its
source. Everything in this round is opt-in and additive.

\section{What the original model does at reversals}
\label{sec:kb-paper-reversals}

A close reading of the source paper \cite{KoBathe2026} (73-page PDF,
read page by page) fixes what is, and what is not, treated:

\begin{itemize}
  \item Loading versus unloading is classified per material point by
        the $J_2$ ratchet $\tau_o \le \tau_o^{\max}$ (Eq.~2b of the
        paper), where $\tau_o^{\max}$ is the maximum over the
        \emph{converged} history, and unloading is solved as an
        incremental \emph{linear} response with softened moduli
        ($K_s$, $G_s$).
  \item Crack closure is a \emph{discrete constitutive branch switch}
        (Table~1 and Eqs.~18--20 of the paper): a crack is open while
        $0 < \bar e^{C}_n \le \max_{\tau}(\bar e^{C}_n)$, and on
        closure the crack count drops ($n \to n-1$) with an
        instantaneous change of constitutive matrix. No transition
        smoothing of any kind is applied.
  \item The celebrated ``nonlinear solution procedure'' of the title
        is the \emph{local} plastic--fracture--crack iteration at the
        material point (Fig.~8 of the paper); globally the model rides
        standard incremental FE solves.
  \item Every structural validation in the paper is \emph{monotonic}
        (panels, beams, benchmark problems). A text-layer search of all
        73 pages finds ``cyclic'' only in the titles of five
        bibliography entries. Structural load reversals, global
        convergence across them, and the selection among the multiple
        equilibria that the discrete closure switch creates are not
        addressed.
\end{itemize}

The reversal pathology characterised in
Caps.~\ref{ch:branch-selection-reversals}--\ref{ch:energy-linesearch-fe2}
(dense spurious basins at reversals, stagnated high-force acceptances,
the need for descent- or control-based selection) is therefore genuinely
outside the scope of the original model, and the regularization
programme of this chapter is new ground.

\section{Why kinematic control regularizes}
\label{sec:kb-control}

The FE\textsuperscript{2} evidence of
Cap.~\ref{ch:energy-linesearch-fe2} showed that once the local history
tracks the macro, the two-way column completes the full protocol with
zero non-converged steps and smooth reversals, while the same material
in the monolithic driver keeps its stagnation band. The mechanical
reason is that the RVE is driven by \emph{imposed section kinematics},
a Dirichlet-dominated control. Every kinematic mode that is imposed
removes one direction of the solution manifold in which spurious
equilibria can live; with enough imposed modes the local problem at
fixed macro state approaches well-posedness even where the material law
is only piecewise smooth. The macro, whose own (fiber or homogenised)
response is smooth, thereby \emph{selects the branch} for the local
model --- the author's ``the macro guides the local solution''
hypothesis, made mechanical.

Until this round the transfer imposed exactly two rigid section planes,
the end faces of the hinge element
(\texttt{LocalBoundaryConditionApplicator}), leaving the whole interior
of the RVE free. The enrichment ladder is
%
\begin{enumerate}
  \item interior section planes at the element-layer levels of the RVE
        mesh (this round, B1),
  \item a 3+-node Timoshenko macro so the hinge element itself carries
        curvature gradients (B2),
  \item the full interpolated boundary field $u(z) = N(\xi)\,u_e$ (B3),
\end{enumerate}
%
complemented by the energy line search of
Cap.~\ref{ch:energy-linesearch-fe2} driving the \emph{local} solve
(D), so that branch selection by descent acts exactly where the
multi-equilibria live.

\section{Fidelity audit of the implementation against the paper}
\label{sec:kb-audit}

A block-by-block comparison of
\path{src/materials/.../KoBatheConcrete3D.hh} against the paper,
carried out with the rendered pages side by side, gives the following
map.

\paragraph{Faithful.} The empirical moduli and coefficients are the
paper's Appendix~A verbatim ($K_e = 11000+3.2f_c$, $G_e$, and the full
$A,b,c,d,k,l,m,n$ sets including both $f_c$ branches), the softened
moduli follow Eq.~10 including the tangential variants
($K_c, G_c$) the paper prescribes for the constitutive matrix, the
mean-stress classification and the Lode angle are standard, and the
loading ratchets ($\tau_o^{\max}$, crack-strain maxima) update on
committed states only, as Table~1 requires.

\paragraph{Deliberate extensions (documented in the header).} Secant
damage-style unloading (the softened moduli are evaluated at the
\emph{historical maxima} $\sigma_o^{\max}, \tau_o^{\max}$ instead of the
paper's last-converged values, making fracture softening irreversible
and unload/reload hysteresis-free), the smoothed tanh crack closure of
the \texttt{production\_stabilized} profile against the paper's
discrete branch switch, the graduated crack shear retention, the
$0.01\,K_e$ floors on the softened moduli, and the delay-damage opening
cap. Each exists to tame precisely the reversal pathologies the paper
never exercises.

\paragraph{One major divergence, not an extension.} The paper's Eq.~2b
defines the classification stress as the \emph{octahedral} shear
stress, $\tau_o = \sqrt{2J_2/3}$ (verified at 300\,dpi, the radical
covers the $/3$; the second form
$\sqrt{\sigma^{dev}_{ij}\sigma^{dev}_{ij}/3}$ confirms it). The
implementation computes
%
\[
  \tau_o^{\text{code}} = \frac{\sqrt{2J_2}}{3}
  = \frac{1}{\sqrt 3}\,\tau_o^{\text{paper}},
\]
%
(both in the 3D header and in the 2D model it inherits from), while
using the paper's Appendix-A constants unchanged. The ratchet
comparisons are scale-invariant, so loading classification is
unaffected, and this is why the model behaves self-consistently. Three
consumers do feel the factor, all in the direction of \emph{less}
degradation than the paper. The shear softening
$G_s = G_e/[1+c(\tau_o/f_c)^{d-1}]$ sees its argument reduced to
$(1/\sqrt3)^{d-1} \approx 0.40$ of the paper's at $f_c=30$; the
internal confining pressure $\sigma_{id}$ is understated by
$1/\sqrt3$; and the cracking envelope of Eq.~16
($\tau_o - 0.944 f_c(\cdot)^{0.724}$, a Kotsovos octahedral fit)
triggers nucleation later. The effect is equivalent to running the
paper's model with rescaled constants
($c' = c\,3^{-(d-1)/2}$, $k' = k\,3^{-n/2}$, envelope inflated by
$\sqrt3$), a stiffer-in-shear, later-cracking variant. Every fall\_n
validation gate to date is self-referential against its own baselines,
so nothing already claimed is invalidated, and correcting the
definition would break bit-identity with every archived campaign.
The correction is therefore an \emph{author decision}: fix and
re-baseline (physically closer to the paper, likely more shear-brittle
and possibly a different reversal landscape), or keep and document the
implicit recalibration. Until that decision, the divergence stands
recorded here.

\section{B1 --- interior-section kinematic transfer}
\label{sec:kb-b1}

The transfer enrichment imposes the beam-interpolated section planes at
the element-layer levels of the RVE mesh ($z_k = k H_{sub}/n_z$,
$\xi_k = -1 + 2k/n_z$; only element boundaries, never the quadratic
mid-levels, so localization can still develop inside each 0.1\,m
layer). The machinery reuses \texttt{extract\_section\_kinematics} at
arbitrary $\xi$ and \texttt{compute\_boundary\_displacements} per
level; \texttt{MultiscaleAnalysis} takes the profile path only when
the local model opts in (concept-detected, building drivers
untouched), and the driver exposes it as
\texttt{--kin-transfer layers}.

The effect was immediate and twofold. The preload steps converge in
\emph{strict} two-way with staggered residuals of
$10^{-5}$--$7\times10^{-4}$ (faces: $10^{-2}$), and the
initial-cracking transient collapses from the $\pm 64$\,kN sign-flip
chaos of the faces run to a single $+26$\,kN release followed by a
clean ramp. It also exposed the last defect in the chain: the first
reversal aborted because a staggered macro failure returned
\texttt{MacroSolveFailed} \emph{without ever reaching} the hybrid
recovery window (so neither the fresh-feedback retry nor the
cleared-injection fallback could act; the frozen $+50$\,mm affine
operator is a barrier no step size crosses, PETSc reason $-11$ at
every bisection depth). Redirecting the staggered failure into the
existing recovery path fixed the reversal, and the unloading branch
now traces the elastic slope cleanly.

\section{D --- energy line search inside the RVE solve}
\label{sec:kb-d}

The evolver gained an opt-in local engine
(\texttt{--local-engine energy-lm}) that drives each kinematic-ramp
sub-increment with the LM continuation over the energy backend
(\texttt{energy\_linesearch} active), reusing the exact machinery
validated on the monolithic driver, with the evolver's
commit discipline unchanged. In the smoke window the trajectories of
\texttt{snes} and \texttt{energy-lm} coincide to four decimals, as
they must where the local solve has no basin ambiguity; the engine's
verdict belongs to the full-protocol campaigns at the large-amplitude
reversals.

\section{Campaign evidence}
\label{sec:kb-evidence}

\paragraph{The corrected physics is harder, and it reorders the
remedies.} The overnight campaigns on the fix branch
(\path{data/output} of the worktree, baseline \path{g0_tauoct}) close
the monolithic side of the two-physics matrix:

\begin{center}
\begin{tabular}{@{}lllll@{}}
\toprule
& \multicolumn{2}{c}{$\tau_o$ transcribed (old)} &
  \multicolumn{2}{c}{$\tau_o$ octahedral (fix)} \\
& noconv & $|V|_{\max}$ & noconv & $|V|_{\max}$ \\
\midrule
baseline          & 80 & 0.051 & 100 & 0.084 \\
line search       & 78 & 0.049 & 97  & 0.059 \\
hybrid TAO        & 77 & 0.038 & 96  & \textbf{0.123} \\
combo             & 77 & 0.055 & 92  & 0.058 \\
CA-tuned (f=1.0)  & 66 & 0.048 & 91  & 0.067 \\
\bottomrule
\end{tabular}
\end{center}

Three conclusions. First, the corrected physics is uniformly harder
(noconv $+20$--$25\%$, baseline peak $+66\%$), consistent with
$40\%$ more effective shear softening and an uninflated cracking
envelope, and its dominant anomaly moves to the \emph{loading} branch
of the $+200$\,mm excursion (steps 154--157). Second, the energy line
search is the only remedy whose character survives the physics change,
containing the peak on both ($0.059$ vs $0.084$) at near-zero cost;
the combo inherits that containment. Third, the hybrid TAO
\emph{inverts}: its trust region, which tamed the old physics' spikes,
commits a $0.123$\,MN state on the corrected one --- the sharper
nonsmoothness pushes BNTR past its useful regime, and TAO should not
be promoted without the line-search guard in front. The CA tuner keeps
finding useful genomes on both physics (and the latched-gene sweep
prefers \texttt{frac}=0.5 with noconv 66 on the corrected physics,
at the price of a different material response --- a physics choice,
not a solver one).

\paragraph{FE\textsuperscript{2} with interior layers --- a genuine
trade-off, honestly measured.} Both full-protocol layer campaigns
complete all 196 steps with zero non-converged rows and \emph{49
strictly converged coupling steps against 20 with end faces}; the
initial-cracking transient disappears, and the sampled reversal
windows (20--22, 116--118, 180--182) unload monotonically along the
elastic slope. The \emph{global} trajectory, however, is rougher than
the end-face run: 21 steps exceed $0.03$\,MN against 5 (12 vs.\ 2
above $0.05$), total variation of $V$ 1.39 vs.\ 1.15\,MN, peak
$0.10$\,MN on late unloading branches. The mechanics is legible:
rigid section planes every $0.1$\,m over-constrain the RVE exactly
where cracking wants to localize between planes, its condensed
response stiffens and spikes, and that excess feeds the macro through
the (now fresh) feedback. Strong kinematic control always selects
\emph{a} branch --- occasionally an artificially stiff one. The
natural refinements are partial layer control (impose only the
lateral translation per level, leaving in-plane warping and axial
free), fewer interior planes, or weak (penalty) enforcement, all of
which live in the same seam this round built.

\paragraph{Startup strategies --- when to engage matters more than
how.} The remaining end-face imperfection was the initial-cracking
transient, and a four-way experiment
(\texttt{--two-way-warmup}~/~\texttt{--coupling-start-step}) settled
it decisively. Immediate coupling gives the known transient (peak
$0.064$\,MN, roughness 1.15). A short one-way warmup through first
cracking (8 steps) already removes it (peak $0.029$, zero steps above
$0.03$\,MN). A \emph{long} warmup through the first push (20 steps) is
the \emph{worst} strategy of all (peak $0.128$): the RVE history
developed under one-way tracking is not consistent with the coupled
equilibrium, and its first feedback jolts the macro --- tracked
history is not consistent history. And the pure delayed start that
engages a \emph{virgin} RVE exactly at the first reversal wins on
every metric (peak $0.027$\,MN, zero steps above $0.03$, roughness
$0.60$ --- smoother than the uncoupled fiber itself, while genuinely
carrying the RVE feedback at the large-amplitude peaks): the
unloading branch is elastic and single-valued, so consistency at
engagement is free, and the damage history then accumulates
progressively inside the limiter-protected loop. The engagement
\emph{point}, not the preparation, is the lever --- reversals are the
natural coupling-activation events, which composes directly with the
criterion-based activation of the building drivers.

The energy-lm local engine produces a trajectory \emph{identical to
four decimals} to the SNES engine across the entire protocol,
excursions included. That is the cleanest confirmation of the
chapter's thesis available: under layer-level kinematic control the
RVE's local solves never face basin ambiguity, so there is nothing
left for the energy criterion to select --- and the residual global
roughness is therefore a \emph{coupling-side} phenomenon (the
feedback of an over-constrained RVE), not a local branch jump. The
local line search remains the right tool where the control is weaker
(end-face transfer, larger sub-models), and it costs nothing to keep
gated.
